\section{Geometry}

In triangle notation, $A,B,C$ denote vertices and $a,b,c$ the opposite side lengths. Lowercase Greek letters such as $\ell,m,n$ may denote lines. The notation $\overrightarrow{AB}$ may denote either the ray from $A$ through $B$ or the vector from $A$ to $B$, according to context. In differential geometry, the Frenet frame follows the standard uppercase convention $\mathbf T,\mathbf N,\mathbf B$, typeset in bold to emphasize that these are vectors; uppercase italic letters such as $X,Y$ may denote vector fields in the usual field-specific convention.

\begin{notationtable}
$\triangle ABC$ & Triangle with vertices $A,B,C$. \\
$\angle ABC$ & Geometric angle with vertex $B$. \\
$m\angle ABC$ & Measure of the angle $\angle ABC$. \\
$\overline{AB}$ & Segment joining $A$ and $B$. \\
$AB$ & Length of $\overline{AB}$. \\
$\overrightarrow{AB}$ & Ray from $A$ through $B$, or vector from $A$ to $B$, according to context. \\
$\overrightarrow{AB}=\mathbf b-\mathbf a$ & Vector from $A$ to $B$ when $\mathbf a$ and $\mathbf b$ are the corresponding position vectors. \\
$\overleftrightarrow{AB}$ & Line through $A$ and $B$. \\
$\ell\parallel m$ & Lines $\ell$ and $m$ are parallel. \\
$\ell\perp m$ & Lines $\ell$ and $m$ are perpendicular. \\
$\triangle ABC\cong\triangle DEF$ & Congruent triangles; corresponding vertices are listed in the same order. \\
$\triangle ABC\sim\triangle DEF$ & Similar triangles; corresponding vertices are listed in the same order. \\
$[ABC]$ & Area of $\triangle ABC$. \\
$C(O,r)$ & Circle with center $O$ and radius $r$. \\
$\odot O$ & Circle centered at $O$, especially in elementary geometry when the radius is understood. \\
$\widehat{AB},\ \widehat{ACB}$ & Arcs specified by their endpoints, with an additional point used when needed to distinguish the intended arc. \\
$O,I,H,G$ & Circumcenter, incenter, orthocenter, and centroid of a triangle, respectively. \\
$R,r$ & Circumradius and inradius of a triangle, respectively. \\
$\displaystyle s=\frac{a+b+c}{2}$ & Semiperimeter of a triangle with side lengths $a,b,c$. \\
\addlinespace[0.35em]
$d(P,Q)$ & Distance between the points $P$ and $Q$. \\
$\dist(P,S)$ & Distance from the point $P$ to the set $S$. \\
$\lengthop(C),\ \ell(C)$ & Length of $C$. \\
$\area(S)$ & Area of $S$. \\
$\vol(S)$ & Volume of $S$. \\
$\diamop(S)$ & Diameter of $S$. \\
$\conv(S)$ & Convex hull of $S$. \\
$\aff(S)$ & Affine hull of $S$. \\
$\mathbf x=\mathbf p+t\mathbf v$ & Parametric equation of a line through $\mathbf p$ in the direction $\mathbf v$. \\
$\mathbf n\cdot(\mathbf x-\mathbf p)=0$ & Hyperplane through $\mathbf p$ with normal vector $\mathbf n$. \\
$\R^n$ & $n$-dimensional real Euclidean space. \\
$S^n=\{\mathbf x\in\R^{n+1}\mid\|\mathbf x\|=1\}$ & Unit $n$-sphere. \\
$S_r^{n-1}(\mathbf a)$ & Sphere of radius $r$ centered at $\mathbf a\in\R^n$. \\
$\Isom(X)$ & Isometry group of the metric or geometric space $X$. \\
\addlinespace[0.35em]
$\mathbb{RP}^n$ & Real projective $n$-space. \\
$\mathbb{CP}^n$ & Complex projective $n$-space. \\
$[x_0:\cdots:x_n]$ & Homogeneous coordinates in projective space. \\
$\mathbb{H}^n$ & $n$-dimensional hyperbolic space. \\
$d_{\mathbb H}$ & Hyperbolic distance, when the hyperbolic metric is understood. \\
\addlinespace[0.35em]
$\gamma\colon I\to\R^n$ & Parametrized curve, where $I\subseteq\R$ is an interval. \\
$s$ & Arc-length parameter, when so specified. \\
$\mathbf T,\ \mathbf N,\ \mathbf B$ & Unit tangent, principal normal, and binormal vectors of the Frenet frame. \\
$\kappa$ & Curvature of a regular curve. \\
$\tau$ & Torsion of a space curve. \\
$T_pM$ & Tangent space of the manifold or regular surface $M$ at $p$. \\
$\mathbf n\colon M\to S^2$ & Gauss map of an oriented surface $M\subseteq\R^3$, assigning its unit normal vector. \\
$\mathrm{I},\ \mathrm{II}$ & First and second fundamental forms of a regular surface. \\
$\kappa_1,\ \kappa_2$ & Principal curvatures of a regular surface. \\
$K=\kappa_1\kappa_2$ & Gaussian curvature of a regular surface. \\
$\displaystyle H=\frac{\kappa_1+\kappa_2}{2}$ & Mean curvature of a regular surface under the average convention. \\
$(M,g)$ & Riemannian manifold $M$ with Riemannian metric $g$. \\
$g_p$ & Riemannian inner product on $T_pM$. \\
$\nabla_XY$ & Covariant derivative of the vector field $Y$ in the direction of the vector field $X$; in Riemannian geometry, $\nabla$ commonly denotes the Levi--Civita connection. \\
$\Gamma^k_{ij}$ & Christoffel symbols of a connection in local coordinates. \\
$\exp_p\colon T_pM\to M$ & Riemannian exponential map at $p$, on its domain of definition. \\
$\nabla_{\dot\gamma}\dot\gamma=0$ & Geodesic equation for an affinely parametrized curve $\gamma$. \\
\end{notationtable}
